% Pass 1335: Brauer-tree completion of the 58|23 bridge.
\section{Brauer-tree completion of the five-primary frame extension}
\label{sec:pass1335-brauer-tree}

Let \(k=\mathbf F_5\), let \(G=W(E_6)\cong U_4(2).2\), and let
\(G'=PSp(4,3)\cong U_4(2)\). GAP/CTblLib gives cyclic Sylow \(5\)-subgroups
of order \(5\). For both groups, the defect-one block containing the
ordinary \(81\)-character has the following shape; for the outer group there
are two sign-twist-related \(81\)-blocks, and GAP checks both:
\[
  1\;-\;23\;-\;58\;-\;6,
\]
where the four entries are the dimensions of the simple edge modules, not the
ordinary vertices.
Equivalently, the ordinary vertices have degrees
\[
  1\;-\;24\;-\;81\;-\;64\;-\;6,
\]
and the \(23\)- and \(58\)-edges meet at the ordinary \(81\)-vertex. The
exceptional multiplicity is one. Brauer-tree algebra theory therefore gives
\[
 \dim_k\operatorname{Ext}^1_G(23,58)
 =\dim_k\operatorname{Ext}^1_G(58,23)=1,
\]
with the same equalities for \(G'\).

Pass~1147 already constructs a nonsplit sequence
\[
 0\longrightarrow I_{58}\longrightarrow S_{81}
 \longrightarrow Q_{23}\longrightarrow0
\]
over both groups. Its class consequently spans the full
\(\operatorname{Ext}^1_G(Q_{23},I_{58})\) and
\(\operatorname{Ext}^1_{G'}(Q_{23},I_{58})\): the nonsplit middle-module
type is unique up to rescaling its two simple endpoints.

The literal \(432\)-character contributes \(2\cdot6+2\cdot64+81\) in this
defect block, so the corresponding rational Hecke-commutant contribution has
dimension \(2^2+2^2+1^2=9\). This is the nonsemisimple nine-dimensional
central block of Pass~1330. The other nine-dimensional block is the
defect-zero species-\(20\) contribution \(M_3(\mathbf F_5)\). A direct GAP
derivation calculation on the certified \(26^3\) multiplication tensor gives
the scalar-simple Ext quiver
\[
 h_6\rightleftarrows h_5\rightleftarrows h_7,
\]
each displayed Ext space being one-dimensional. Its radical powers have
dimensions \(6,2,0\). Here \(h_i\) is the \(i\)-th one-dimensional simple
\(\mathcal H_{26,\mathbf F_5}\)-module in the frozen seven-character order;
the subscripts are Hecke-character indices, not module dimensions.

This Hecke block is a condensation shadow of the same cyclic-defect group
block, not the literal \(58|23\) middle module. The full group Brauer tree,
not the Hecke quiver alone, proves the Ext theorem.

Finally,
\[
 \mathbf F_5[C_3]\cong\mathbf F_5\times\mathbf F_{25}.
\]
Since characteristic \(5\) divides neither \(|C_3|\) nor \(|S_3|\), color
triality is semisimple and cross-color Ext vanishes. The colored middle
space splits as \(81+162\) over \(\mathbf F_5\), and as three independent
\(81\)-extensions after scalar extension to \(\mathbf F_{25}\). Triality
transports the unique class; it neither creates it nor selects a physical
copy.
